% ============================================================================
% MANUSCRITO BORRADOR - Entrega 3 (2A)
% Titulo: Comparacion de la calidad de requisitos funcionales elicitados por
%         analistas humanos frente a los generados por un modelo de lenguaje
%         (LLM) en un sistema agricola con IA
% Sistema AgriMoreira - Ingenieria de Requerimientos [20303] - UTEQ
% Integrante responsable: Escudero Plaza Maria del Rosario
% Nota: secciones producidas en la Entrega 3 (2A): introduccion, trabajo
% relacionado y metodologia. Resultados, discusion, amenazas a la validez y
% conclusiones se completan en la Entrega 4 (2B) con datos reales.
% La plantilla oficial de la revista objetivo (Springer) se aplica al envio.
% ============================================================================
\documentclass[11pt,a4paper]{article}

\usepackage[utf8]{inputenc}
\usepackage[T1]{fontenc}
\usepackage[spanish,es-tabla]{babel}
\usepackage[margin=2.5cm]{geometry}
\usepackage{amsmath,amssymb}
\usepackage{booktabs}
\usepackage{graphicx}
\usepackage{hyperref}
\usepackage{longtable}
\usepackage{array}
\usepackage{multirow}
\usepackage{xcolor}

\hypersetup{colorlinks=true,linkcolor=blue,urlcolor=blue,citecolor=blue}

\title{\textbf{Comparacion de la calidad de requisitos funcionales elicitados
por analistas humanos frente a los generados por un modelo de lenguaje (LLM)
en un sistema agricola con IA}}
\author{Escudero Plaza Maria del Rosario\\[0.3em]
\small Universidad Tecnica Estatal de Quevedo (UTEQ)\\
\small Proyecto Fin de Curso - Ingenieria de Requerimientos [20303]}
\date{Version 1.0 (borrador, Entrega 3) \textbullet\ Semestre 2026-2027}

\begin{document}

\maketitle

% ============================================================================
\section*{Resumen}
% ============================================================================
Los modelos de lenguaje de gran escala (LLM) se estan integrando cada vez mas
en la fase de elicitacion de requisitos de la ingenieria de software
\cite{cheng2026,arora2024}. Sin embargo, la evidencia empirica sobre la
calidad de los requisitos funcionales (RF) generados por un LLM frente a los
elicitados por analistas humanos a partir de un mismo material fuente es
todavia limitada, especialmente en dominios agroindustriales de paises en
desarrollo. En este trabajo se realiza un cuasi-experimento apareado en el
dominio agroindustrial ecuatoriano: un equipo humano y un LLM (gpt-4o-2024-08-06,
temperatura 0.0, top-p 1.0, semilla 42) producen sendos conjuntos de
requisitos funcionales (minimo 25 cada uno) a partir de la misma transcripcion
anonimizada de una entrevista de campo del sistema \textbf{AgriMoreira}. Tres o
mas evaluadores expertos independientes puntuan a ciegas ambos conjuntos en
cinco dimensiones de calidad (completitud, ausencia de ambiguedad,
verificabilidad, correccion respecto de la fuente y consistencia interna). El
acuerdo inter-evaluador se mide con kappa de Cohen \cite{cohen1960} y de
Fleiss \cite{fleiss1971}; las medias se comparan con prueba t pareada o
Wilcoxon segun normalidad, reportando el tamano del efecto (d de Cohen o delta
de Cliff) y el calculo de potencia previo (alpha = 0.05, 1-beta = 0.80)
\cite{molleri2020}. El protocolo se registra de forma previa en el OSF
\cite{nosek2018} y el conjunto de datos anonimizado se depositara en Zenodo
bajo licencia CC BY 4.0 \cite{wilkinson2016}.

\noindent\textbf{Palabras clave:} ingenieria de requisitos, elicitacion,
modelos de lenguaje, calidad de requisitos, experimento empirico.

% ============================================================================
\section{Introduccion}
% ============================================================================

El sistema \textbf{AgriMoreira} es una plataforma de gestion agricola con
inteligencia artificial para cultivos de platanero verde y cacao en Manabi,
Ecuador, desarrollada como proyecto academico de la asignatura de Ingenieria
de Requerimientos. La finca real de la organizacion cliente (Agricola
Moreira, seudonimo declarado en el protocolo de anonimizacion) registra hoy
su produccion, inventario y actividades de forma manual en cuadernos, segun
las entrevistas y encuestas de la primera y segunda ronda de campo. En este
dominio, la calidad de los requisitos funcionales que sustentan el sistema es
determinante: un requisito ambiguo o incompleto se traduce directamente en
perdida de informacion y en decisiones agronomicas incorrectas. Por ello,
estudiar la ingenieria de requerimientos en este contexto no es un ejercicio
academico abstracto, sino una necesidad operativa real de una cadena de valor
agropecuaria.

El gap de conocimiento que este trabajo aborda es doble. Por un lado, la
revision sistematica de Cheng et al.\ \cite{cheng2026} indica que la fase de
gestion de requisitos recibe atencion insuficiente frente a otras actividades
del proceso, y que las guias de uso de LLM en ingenieria de requisitos, como
la de Vogelsang y Fischbach \cite{vogelsang2024}, aun no han sido validadas en
contextos latinoamericanos ni en dominios agropecuarios. Por otro lado,
Arora et al.\ \cite{arora2024} senalan que, aunque los LLM pueden acelerar la
elicitacion, se carece de evidencia empirica controlada que compare la calidad
resultante frente a la elicitacion humana clasica sobre una misma fuente.

Las contribuciones de este trabajo son las siguientes:
\begin{enumerate}
    \item Una evidencia empirica adicional que complementa el mapeo
    sistematico de Montgomery et al.\ \cite{montgomery2022} y las guias de
    \cite{vogelsang2024} sobre el uso de LLM en la elicitacion de requisitos,
    medida en cinco dimensiones de calidad con evaluadores a ciegas.
    \item Un protocolo reproducible y preregistrado (OSF) para comparar RF
    humanos frente a LLM, con instrumentos, consignas, parametros del modelo y
    scripts de analisis publicados en el repositorio.
    \item Un conjunto de datos anonimizado depositado en Zenodo con licencia
    CC BY 4.0, siguiendo los principios FAIR \cite{wilkinson2016} y las
    practicas de citacion de software \cite{smith2016,druskat2021}.
\end{enumerate}

El resto del manuscrito se organiza asi: la Seccion 2 presenta el trabajo
relacionado; la Seccion 3 describe la metodologia del cuasi-experimento
(Enfoque 1); las secciones de resultados, discusion y conclusiones se
completaran en la Entrega 4 (2B) con los datos analizados.

% ============================================================================
\section{Trabajo relacionado}
% ============================================================================

Para delimitar el estado del arte se ejecuto una busqueda sistematica de
literatura con la cadena de tres terminos: \texttt{("requirements engineering"
OR "requisitos de software") AND ("large language model" OR "ChatGPT" OR
"Claude") AND ("elicitation" OR "quality")}, en las bases de datos Scopus y
Google Scholar, limitando el periodo a los ultimos tres anos. Los criterios de
inclusion fueron: estudios revisados por pares con datos empiricos publicados
en conferencias o revistas de ingenieria de requisitos. Los criterios de
exclusion fueron: editoriales, presentaciones breves de menos de cuatro paginas
y estudios sin datos.

El Cuadro~\ref{tab:relacionado} resume los estudios mas relevantes.

\begin{table}[htbp]
\centering
\small
\begin{tabular}{p{2.6cm} p{0.9cm} p{1.7cm} p{2.2cm} p{3.2cm} p{3.0cm}}
\toprule
\textbf{Referencia} & \textbf{Ano} & \textbf{Tipo de estudio} & \textbf{Poblacion} & \textbf{Intervencion} & \textbf{Diferencia con nuestro trabajo} \\
\midrule
Cheng et al.\ \cite{cheng2026} & 2026 & Revision sistematica & Literatura de IA generativa en RE & Analisis de estudios sobre generacion de requisitos & Es una revision; no mide calidad en un experimento controlado \\
Arora et al.\ \cite{arora2024} & 2024 & Vision/posicion & Comunidad de RE & Evaluacion cualitativa del rol de los LLM & No aporta comparacion cuantitativa apareada \\
Vogelsang y Fischbach \cite{vogelsang2024} & 2024 & Guia & Practicantes de RE & Directrices de uso de LLM en tareas de PLN & No valida la guia en dominio agroindustrial latinoamericano \\
Montgomery et al.\ \cite{montgomery2022} & 2022 & Mapeo sistematico & Estudios de calidad de requisitos & Clasificacion de operacionalizaciones de calidad & No compara humanos frente a LLM sobre una misma fuente \\
\bottomrule
\end{tabular}
\caption{Cuadro comparativo del trabajo relacionado.}
\label{tab:relacionado}
\end{table}

El nicho concreto que ocupa nuestro estudio es la comparacion controlada y
cuantitativa, sobre un unico material fuente, entre la calidad de los RF
producidos por un equipo humano y la producida por un LLM con parametros
fijos, en un dominio agroindustrial real de Ecuador. Los estudios previos
documentan el potencial de los LLM o clasifican las dimensiones de calidad,
pero no ofrecen una evidencia apareada de cinco dimensiones con evaluadores a
ciegas y analisis estadistico preregistrado en este contexto.

% ============================================================================
\section{Metodologia}
% ============================================================================

La metodologia sigue la plantilla de reporte de experimentos de Jedlitschka
et al.\ \cite{jedlitschka2008} y la guia de Wohlin et al.\ \cite{wohlin2012}
para cuasi-experimentos. El detalle completo esta en el protocolo experimental
\cite{protocolo}, registrado en el OSF con fecha anterior a la ejecucion.

\subsection{Pregunta de investigacion (PICOC)}
\label{subsec:rq}
La pregunta principal es: \emph{?`En que dimensiones de calidad los RF
elicitados por un equipo humano difieren de los generados por un LLM a partir
del mismo material fuente?}

\begin{center}
\begin{tabular}{>{\raggedright\arraybackslash}p{3cm} p{12cm}}
\toprule
\textbf{Componente} & \textbf{Descripcion} \\
\midrule
Poblacion (P) & Conjuntos de RF del sistema AgriMoreira (>= 25 RF humanos del
ERS y >= 25 RF generados por LLM). \\
Intervencion (I) & Evaluacion a ciegas con rubrica de cinco dimensiones
(escala 1--5). \\
Comparacion (C) & Conjunto A (LLM) frente a Conjunto B (equipo humano). \\
Resultado (O) & Puntuacion media de calidad por dimension, acuerdo
inter-evaluador (kappa) y tamano del efecto. \\
Contexto (C) & Dominio agroindustrial ecuatoriano, periodo 2026-2027. \\
\bottomrule
\end{tabular}
\end{center}

\subsection{Hipotesis}
Para cada dimension de calidad \(j \in \{1, \ldots, 5\}\):
\[
H_{0j}: \mu_{B,j} = \mu_{A,j}, \qquad
H_{1j}: \mu_{B,j} \neq \mu_{A,j},
\]
con \(\alpha = 0{,}05\) y potencia \(1-\beta = 0{,}80\).

\subsection{Variables}
\begin{itemize}
    \item \textbf{Independiente:} origen del conjunto de requisitos (A = LLM,
    B = humano).
    \item \textbf{Dependientes:} puntuacion de calidad por dimension (1--5) y
    acuerdo inter-evaluador (kappa de Cohen y de Fleiss).
    \item \textbf{De control:} material fuente unico (transcripcion
    anonimizada EV-01), rubrica identica, evaluadores ciegos e independientes,
    y parametros del LLM fijos (modelo, temperatura, top-p y semilla).
\end{itemize}

\subsection{Diseno del estudio}
Cuasi-experimento apareado \cite{wohlin2012}. El material fuente unico es la
transcripcion anonimizada de la entrevista al administrador (Entr-01 = EV-01),
disponible en \texttt{02\_Evidencias/Transcripciones/}. El LLM recibe
exclusivamente esa transcripcion con la consigna de la guia y produce el
Conjunto A; el Conjunto B lo conforman los RF del ERS (RF-01 a RF-16 segun la
matriz de trazabilidad). Cada conjunto debe alcanzar un minimo de 25 requisitos
individuales para el analisis estadistico apareado \cite{montgomery2022}.

\subsection{Participantes y reclutamiento}
Se contratan de 3 a 5 evaluadores expertos independientes (por ejemplo,
estudiantes de otro paralelo o docentes del area de Ingenieria de
Requerimientos) con criterios de inclusion (conocimientos de RE, no haber
participado en la elicitacion) y exclusion (miembros del equipo). Cada
evaluador firma consentimiento informado conforme a la Ley Organica de
Proteccion de Datos Personales del Ecuador \cite{lopdp2021}, con autorizacion
para el uso de datos anonimizados en publicaciones y en el deposito en Zenodo
con CC BY 4.0.

\subsection{Instrumentos}
\begin{enumerate}
    \item Material fuente unico: transcripcion anonimizada de la entrevista
    EV-01.
    \item Consigna exacta del LLM con parametros registrados en
    \texttt{06\_Experimento/prompts\_llm/experimento\_enf1\_prompt.md}.
    \item Rubrica de evaluacion experta (cinco dimensiones, escala 1--5) en
    \texttt{06\_Experimento/instrumentos/}.
    \item Hoja de evaluacion a ciegas
    (\texttt{06\_Experimento/instrumentos/hoja\_evaluacion\_a\_ciegas.csv}).
    \item Plantilla de consentimiento actualizado (carpeta \texttt{08\_Etica/}).
\end{enumerate}

\subsection{Plan de analisis estadistico}
\begin{enumerate}
    \item Normalidad de las diferencias apareadas: prueba de Shapiro-Wilk.
    \item Comparacion de medias: t pareada (normales) o Wilcoxon (no normales).
    \item Tamano del efecto: d de Cohen o delta de Cliff.
    \item Correccion por comparaciones multiples (Bonferroni).
    \item Acuerdo inter-evaluador: kappa de Cohen por pares y kappa de Fleiss.
    \item Potencia estadistica previa: alpha 0.05, 1-beta 0.80
    \cite{molleri2020}.
\end{enumerate}
Todos los calculos se ejecutan con los scripts versionados en
\texttt{06\_Experimento/scripts\_analisis/analisis\_ef1.py} a partir de los
datos crudos de \texttt{06\_Experimento/resultados/}.

\bibliographystyle{unsrt}
\bibliography{referencias}

\end{document}
